% Options for packages loaded elsewhere
\PassOptionsToPackage{unicode}{hyperref}
\PassOptionsToPackage{hyphens}{url}
\documentclass[
]{article}
\usepackage{xcolor}
\usepackage{amsmath,amssymb}
\setcounter{secnumdepth}{-\maxdimen} % remove section numbering
\usepackage{iftex}
\ifPDFTeX
  \usepackage[T1]{fontenc}
  \usepackage[utf8]{inputenc}
  \usepackage{textcomp} % provide euro and other symbols
\else % if luatex or xetex
  \usepackage{unicode-math} % this also loads fontspec
  \defaultfontfeatures{Scale=MatchLowercase}
  \defaultfontfeatures[\rmfamily]{Ligatures=TeX,Scale=1}
\fi
\usepackage{lmodern}
\ifPDFTeX\else
  % xetex/luatex font selection
\fi
% Use upquote if available, for straight quotes in verbatim environments
\IfFileExists{upquote.sty}{\usepackage{upquote}}{}
\IfFileExists{microtype.sty}{% use microtype if available
  \usepackage[]{microtype}
  \UseMicrotypeSet[protrusion]{basicmath} % disable protrusion for tt fonts
}{}
\makeatletter
\@ifundefined{KOMAClassName}{% if non-KOMA class
  \IfFileExists{parskip.sty}{%
    \usepackage{parskip}
  }{% else
    \setlength{\parindent}{0pt}
    \setlength{\parskip}{6pt plus 2pt minus 1pt}}
}{% if KOMA class
  \KOMAoptions{parskip=half}}
\makeatother
\usepackage{graphicx}
\makeatletter
\newsavebox\pandoc@box
\newcommand*\pandocbounded[1]{% scales image to fit in text height/width
  \sbox\pandoc@box{#1}%
  \Gscale@div\@tempa{\textheight}{\dimexpr\ht\pandoc@box+\dp\pandoc@box\relax}%
  \Gscale@div\@tempb{\linewidth}{\wd\pandoc@box}%
  \ifdim\@tempb\p@<\@tempa\p@\let\@tempa\@tempb\fi% select the smaller of both
  \ifdim\@tempa\p@<\p@\scalebox{\@tempa}{\usebox\pandoc@box}%
  \else\usebox{\pandoc@box}%
  \fi%
}
% Set default figure placement to htbp
\def\fps@figure{htbp}
\makeatother
\setlength{\emergencystretch}{3em} % prevent overfull lines
\providecommand{\tightlist}{%
  \setlength{\itemsep}{0pt}\setlength{\parskip}{0pt}}
\usepackage{bookmark}
\IfFileExists{xurl.sty}{\usepackage{xurl}}{} % add URL line breaks if available
\urlstyle{same}
\hypersetup{
  hidelinks,
  pdfcreator={LaTeX via pandoc}}

\author{}
\date{}

\begin{document}

\begin{quote}
\includegraphics[width=2.4224in,height=0.82448in]{media/image1.png}

\textbf{Technologie du commerce électronique}
\end{quote}

\section{INF27523 (MS)}\label{inf27523-ms}

\begin{quote}
UNIVERSIT´E DU QU´EBEC A` RIMOUSKI

D´EPARTEMENT DE MATH´EMATIQUES, D'INFORMATIQUE ET DE G´ENIE
\end{quote}

RAPPORT

\section{Travail Pratique 01}\label{travail-pratique-01}

\begin{quote}
\textbf{E´QUIPE :}
\end{quote}

Amina Missoum

Émile Santerre

Léa Esso-Solam Bikassa

\begin{quote}
\textbf{Chargée de cours :}
\end{quote}

\subsection{Oumeima Nuigues}\label{oumeima-nuigues}

Table des matières

\section{Table des matières}\label{table-des-matiuxe8res}

\hyperref[introduction]{1 Introduction \hyperref[introduction]{3}}

\hyperref[description-du-projet]{2 Description du projet
\hyperref[description-du-projet]{3}}

\hyperref[technologies-utilisuxe9es]{3 Technologies utilisées
\hyperref[technologies-utilisuxe9es]{4}}

\hyperref[mise-en-ux153uvre-et-fonctionnalituxe9s]{4 Mise en œuvre et
fonctionnalités \hyperref[mise-en-ux153uvre-et-fonctionnalituxe9s]{5}}

\hyperref[duxe9fis-rencontruxe9s]{5 Défis rencontrés
\hyperref[duxe9fis-rencontruxe9s]{7}}

\hyperref[conclusion]{6 Conclusion \hyperref[conclusion]{8}}

\section{}\label{section}

\begin{quote}
1 Introduction
\end{quote}

\section{1 Introduction}\label{introduction}

Dans le cadre du cours de commerce électronique, ce travail pratique
avait pour objectif de concevoir et d'implémenter une application web de
boutique en ligne, en mettant en pratique les notions vues en classe
autour des architectures MVC, de la gestion de sessions et de
l'intégration d'une passerelle de paiement. Nous avons développé un site
permettant à un client de naviguer parmi les restaurants, de consulter
leur menu, de constituer un panier, puis de finaliser sa commande à
l'aide d'un paiement sécurisé avec Stripe.

Au-delà de la prise de commande, le projet nous a permis d'aborder
plusieurs aspects clés d'une application de commerce électronique :
gestion de comptes utilisateurs (client et vendeur), association des
restaurants à un vendeur, validation métier (un panier ne peut contenir
que des plats d'un seul restaurant), gestion du panier en session,
historique des commandes et interface dédiée aux restaurateurs pour
consulter leurs ventes. Ce rapport présente la description globale du
projet, les technologies utilisées, la mise en œuvre des principales
fonctionnalités, les défis rencontrés ainsi que les apprentissages que
nous en retirons.

2 Description du projet

\section{Description du projet}\label{description-du-projet}

\subsubsection{}\label{section-1}

\subsubsection{\texorpdfstring{Le projet consiste en la réalisation d'un
site web de commande de repas nommé « ÉléganceShop », inspiré d'une
plateforme de type agrégateur de restaurants, mais centré sur un
ensemble restreint de restaurants. L'application vise à offrir à un
client une expérience fluide allant de la découverte du menu jusqu'au
paiement en ligne, tout en fournissant aux vendeurs un accès simple aux
commandes reçues.\\
}{Le projet consiste en la réalisation d'un site web de commande de repas nommé « ÉléganceShop », inspiré d'une plateforme de type agrégateur de restaurants, mais centré sur un ensemble restreint de restaurants. L'application vise à offrir à un client une expérience fluide allant de la découverte du menu jusqu'au paiement en ligne, tout en fournissant aux vendeurs un accès simple aux commandes reçues. }}\label{le-projet-consiste-en-la-ruxe9alisation-dun-site-web-de-commande-de-repas-nommuxe9-uxe9luxe9ganceshop-inspiruxe9-dune-plateforme-de-type-agruxe9gateur-de-restaurants-mais-centruxe9-sur-un-ensemble-restreint-de-restaurants.-lapplication-vise-uxe0-offrir-uxe0-un-client-une-expuxe9rience-fluide-allant-de-la-duxe9couverte-du-menu-jusquau-paiement-en-ligne-tout-en-fournissant-aux-vendeurs-un-accuxe8s-simple-aux-commandes-reuxe7ues.}

\subsubsection{Plus concrètement, nous avons
voulu~:}\label{plus-concruxe8tement-nous-avons-voulu}

\begin{itemize}
\item
  Proposer une interface claire pour parcourir les restaurants et leurs
  plats.
\item
  Permettre la création de comptes Client et Vendeur, avec un rôle bien
  défini.
\item
  Une liste des jeux de données avec recherche et filtres dynamiques.
\item
  Gérer un panier persistant, lié à la session, avec calcul automatique
  du total.
\item
  Intégrer un paiement réel via Stripe pour simuler un scénario de
  production.
\item
  Enregistrer les commandes et permettre au restaurant associé d'en
  consulter le détail.
\end{itemize}

\subsubsection{}\label{section-2}

2.2 Contexte fonctionnel

D'un point de vue opérationnel, l\textquotesingle application se divise
en deux parcours utilisateurs. Pour le client, le scénario commence par
la création d\textquotesingle un compte ou une authentification
sécurisée. Une fois connecté, l\textquotesingle utilisateur navigue
parmi les restaurants disponibles, consulte les cartes détaillées et
compose son panier. Le processus termine lors de la phase de validation
où le client accède à une interface de paiement pour confirmer sa
transaction avant de pouvoir suivre l\textquotesingle état de sa
commande dans son historique personnel.

À l\textquotesingle opposé, le vendeur bénéficie d\textquotesingle une
interface d\textquotesingle administration. Lors de son inscription,
celui-ci est obligatoirement associé à un établissement existant, ce qui
lui octroie le droit de consulter la liste des commandes reçues. Cette
interface permet au restaurateur d\textquotesingle accéder aux détails
de chaque vente, tels que la liste des plats commandés, les quantités
exactes, les montants perçus ainsi que les informations de contact du
client.

2.3 Fonctionnalités principales

La mise en œuvre technique de l\textquotesingle application repose sur
plusieurs piliers fonctionnels interconnectés. La gestion des comptes
constitue la première brique, utilisant une authentification par cookies
et des rôles bien définis (Client ou Vendeur) pour restreindre
l\textquotesingle accès aux fonctionnalités spécifiques.
L\textquotesingle exploration des restaurants a été optimisée par la
création de pages détaillées affichant les menus, les descriptions de
chaque plat.

Le cœur de l\textquotesingle application réside cependant dans la
logique du panier d\textquotesingle achat. Ce dernier
n\textquotesingle est pas seulement un espace de stockage temporaire, il
intègre une règle stricte qui limite le contenu à un seul restaurant par
commande pour garantir la cohérence logistique. Cette contrainte est
gérée par un système de détection de conflit qui alerte
l\textquotesingle utilisateur en cas d\textquotesingle ajout incohérent.
Enfin, le processus de paiement est sécurisé par Stripe Elements,
assurant que les données bancaires ne transitent jamais en clair sur
notre serveur. Une fois le paiement validé, le système génère
automatiquement une commande et une facture, mettant ainsi à jour
l\textquotesingle historique du client et le tableau de bord des ventes
du vendeur.

3 Technologies utilisées

\section{3 Technologies utilisées}\label{technologies-utilisuxe9es}

\subsubsection{}\label{section-3}

3.1 Environnement de développement

\subsubsection{}\label{section-4}

Pour ce projet nous nous sommes appuyés sur un écosystème
d\textquotesingle outils professionnels. Le développement a été
principalement orchestré sous Visual Studio, permettant une gestion
fluide du code C\#, des vues Razor et des feuilles de style CSS. Pour
assurer la durabilité du travail et faciliter les phases
d\textquotesingle intégration, nous avons utilisé Git avec GitHub comme
système de gestion de versions. Cette approche nous a permis de
travailler sur des branches distinctes, notamment pour isoler les tests
de l\textquotesingle API Stripe avant de fusionner les fonctionnalités
validées dans la branche principale. L\textquotesingle aspect technique
repose sur le framework ASP.NET Core en MVC, qui structure
l\textquotesingle application selon un modèle de séparation des
responsabilités, tandis qu\textquotesingle Entity Framework Core assure
la liaison entre le code et la base de données relationnelle.

\subsubsection{}\label{section-5}

3.2 Langages et frameworks

Notre application est propulsée par le langage C\# via le framework
ASP.NET Core, qui gère l\textquotesingle ensemble de la logique du
serveur, des contrôleurs spécialisés aux mécanismes de routage. Pour la
partie interface, nous avons privilégié Razor, facilitant la génération
dynamique de contenu HTML à partir des données du modèle, comme
l\textquotesingle affichage des paniers ou des historiques de commandes.

La persistance des données est opérée par Entity Framework Core,
agissant comme une couche d\textquotesingle abstraction pour manipuler
les entités telles que les utilisateurs, les menus et les factures sans
injection SQL directe. Le volet transactionnel est quant à lui sécurisé
par Stripe .NET et Stripe.js, offrant une gestion des intentions de
paiement. Enfin, l\textquotesingle aspect visuel et
l\textquotesingle ergonomie ont été standardisés grâce à Bootstrap 5,
complété par Bootstrap Icons pour une navigation intuitive, le tout
soutenu par une charte graphique personnalisée en HTML5 et CSS3.

3.3 Structure et services

Notre architecture interne a été pensée pour favoriser la modularité et
la réutilisation du code. Nous avons notamment mis en place un service
dédié, le CartService, dont la responsabilité unique est
d\textquotesingle orchestrer la logique du panier
d\textquotesingle achat de sa création initiale en session
jusqu\textquotesingle à la gestion des incrémentations de quantités et
de la suppression d\textquotesingle items.

Ce service interagit avec un modèle de données structuré où chaque
entité possède des relations claires. Un restaurant rassemble plusieurs
plats et se lie à un vendeur unique, tandis qu\textquotesingle un panier
regroupe des éléments référençant des produits spécifiques. Cette
organisation rigoureuse entre les entités User, Restaurant, Plat et
Order permet de maintenir une intégrité des données, particulièrement
lors de la transformation finale d\textquotesingle un panier en commande
après validation du paiement.

4 Mise en œuvre et fonctionnalités

\section{4 Mise en œuvre et
fonctionnalités}\label{mise-en-ux153uvre-et-fonctionnalituxe9s}

\subsubsection{}\label{section-6}

4.1 Organisation des contrôleurs et des vues

L'architecture de l'application a été segmentée en plusieurs pôles de
responsabilité afin de garantir une maintenance aisée et une séparation
claire des préoccupations. Le cœur de l'expérience utilisateur est
piloté par le HomeController et le RestaurantsController, qui gèrent
respectivement la vitrine de la plateforme et
l\textquotesingle affichage des menus. La gestion sensible des accès et
des comptes est isolée dans l\textquotesingle AccountController, tandis
que le cycle de vie de la commande est réparti entre le CartController
pour la phase de sélection, le CheckoutController pour la transaction
financière, et les contrôleurs OrdersController et SalesController pour
le suivi post-achat. Chaque contrôleur communique avec des vues Razor
spécifiques, structurées pour offrir une interface cohérente.
L\textquotesingle utilisation d\textquotesingle un fichier de style
global assure que les composants visuels, tels que les cartes de
restaurants ou les résumés de commandes, conservent une unité esthétique
sur l\textquotesingle ensemble du site.

4.2 Authentification, rôles et profil

La sécurité et la personnalisation de l\textquotesingle expérience
reposent sur un système d\textquotesingle authentification basé sur les
identités stockées en base de données. Lors de la phase
d\textquotesingle inscription, l\textquotesingle utilisateur se voit
assigner un rôle spécifique (Client ou Vendeur) qui détermine ses droits
d\textquotesingle accès au sur de la plateforme. Le processus de
connexion utilise des cookies d\textquotesingle authentification pour
maintenir la session active, permettent de vérifier instantanément le
rôle de l\textquotesingle utilisateur lors de chaque requête. Cette
structure permet, par exemple, de restreindre l\textquotesingle accès au
formulaire de paiement aux seuls clients, ou de réserver les suivis de
vente au restaurateur approprié. Une page de profil dédiée permet aussi
à chaque utilisateur de gérer ses informations personnelles,
garantissant ainsi la conformité aux standards de gestion de comptes
modernes.

4.3 Gestion du panier et règle métier

\subsubsection{}\label{section-7}

\subsubsection{Le panier d'achat constitue l'élément le plus complexe de
la logique métier. Pour éviter toute confusion lors de la préparation
des commandes, nous avons instauré une contrainte technique interdisant
le mélange de plats provenant de différents établissements. Cette règle
est pilotée par le CartService qui, lors de chaque ajout, compare
l\textquotesingle identifiant du restaurant du nouveau plat avec celui
des articles déjà présents. Si une incohérence est détectée, le système
lève une exception qui est interceptée par le contrôleur pour déclencher
une interface modale. Cette fenêtre contextuelle offre au client une
alternative claire, abandonner l\textquotesingle ajout pour préserver
ses choix actuels, ou réinitialiser complètement son panier pour entamer
une nouvelle commande auprès du second restaurant. Cette approche
garantit que chaque commande générée est dirigée vers un seul et unique
restaurant.}\label{le-panier-dachat-constitue-luxe9luxe9ment-le-plus-complexe-de-la-logique-muxe9tier.-pour-uxe9viter-toute-confusion-lors-de-la-pruxe9paration-des-commandes-nous-avons-instauruxe9-une-contrainte-technique-interdisant-le-muxe9lange-de-plats-provenant-de-diffuxe9rents-uxe9tablissements.-cette-ruxe8gle-est-pilotuxe9e-par-le-cartservice-qui-lors-de-chaque-ajout-compare-lidentifiant-du-restaurant-du-nouveau-plat-avec-celui-des-articles-duxe9juxe0-pruxe9sents.-si-une-incohuxe9rence-est-duxe9tectuxe9e-le-systuxe8me-luxe8ve-une-exception-qui-est-interceptuxe9e-par-le-contruxf4leur-pour-duxe9clencher-une-interface-modale.-cette-fenuxeatre-contextuelle-offre-au-client-une-alternative-claire-abandonner-lajout-pour-pruxe9server-ses-choix-actuels-ou-ruxe9initialiser-compluxe8tement-son-panier-pour-entamer-une-nouvelle-commande-aupruxe8s-du-second-restaurant.-cette-approche-garantit-que-chaque-commande-guxe9nuxe9ruxe9e-est-diriguxe9e-vers-un-seul-et-unique-restaurant.}

\subsubsection{}\label{section-8}

4.4 Intégration de Stripe

Le déploiement de la solution de paiement s'est construit autour
d\textquotesingle une communication bidirectionnelle entre notre serveur
et les API de Stripe. La procédure débute par la création
d\textquotesingle une intention de paiement (PaymentIntent) dès que le
client accède à la page de validation, calculant le montant exact en
cents pour éviter les erreurs d\textquotesingle arrondi. Côté client,
nous utilisons Stripe Elements pour injecter un champ de saisie de carte
bancaire sécurisé, ce qui permet de traiter les informations sensibles
sans qu\textquotesingle elles ne transitent par nos propres serveurs,
assurant ainsi une conformité PCI. Une fois que Stripe confirme la
validité de la transaction, un identifiant unique est renvoyé à notre
contrôleur. Ce dernier effectue une vérification finale de sécurité côté
serveur avant de confirmer la création de la commande et de vider le
panier, transformant ainsi une intention d\textquotesingle achat en une
transaction finalisée et enregistrée.

5 Défis rencontrés

\section{5 Défis rencontrés}\label{duxe9fis-rencontruxe9s}

\section{}\label{section-9}

5.1 Gestion du panier et résolution des conflits de restaurants

L'un des défis techniques les plus important a été de maintenir la
cohérence du panier face aux comportements imprévus des utilisateurs.
Initialement, le système permettait l\textquotesingle ajout de plats
issus de différents établissements, ce qui créait une ambiguïté majeure
lors de la répartition des paiements et de la notification aux
restaurateurs. Pour remédier à cela, nous avons conçu une logique
d\textquotesingle exception personnalisée au sein du CartService. La
difficulté ne résidait pas seulement dans l\textquotesingle interdiction
de l\textquotesingle ajout, mais dans la gestion de
l\textquotesingle expérience utilisateur qui en découle. Il a fallu
orchestrer une communication entre le service, le contrôleur et une vue
partielle asynchrone pour que la modale de conflit puisse offrir une
option de réinitialisation automatique du panier. Ce mécanisme a
nécessité une réflexion approfondie sur la gestion des états de session
et sur la persistance des données temporaires.

5.2 Intégration de la passerelle Stripe et synchronisation

L'intégration d\textquotesingle une solution de paiement tierce comme
Stripe a introduit une complexité supplémentaire liée à la
synchronisation entre le front-end et le back-end. L\textquotesingle un
des obstacles majeurs a été de garantir que le cycle de vie du
PaymentIntent corresponde exactement aux actions de
l\textquotesingle utilisateur. Nous avons dû faire face à des situations
où le script Stripe.js ne se chargeait pas correctement ou dont le rendu
visuel était altéré par des contraintes CSS, rendant le champ de saisie
invisible ou inutilisable. De plus, la sécurisation du processus a exigé
une validation rigoureuse côté serveur, il était essentiel de confirmer
le statut du paiement directement via l\textquotesingle API Stripe avant
de valider la commande en base de données, afin d\textquotesingle éviter
la création de commandes "fantômes" en cas
d\textquotesingle interruption de la connexion ou de fraude potentielle.

5.3 Sécurisation des accès et gestion des rôles

La mise en œuvre d\textquotesingle un système à deux types
d\textquotesingle utilisateurs (Client et Vendeur) a apporté des enjeux
de sécurité importants. Le défi consistait à restreindre
l\textquotesingle accès aux URLs sensibles pour éviter
qu\textquotesingle un client ne puisse visualiser les ventes
d\textquotesingle un restaurant. Nous avons rencontré des difficultés
initiales avec la mise à jour des "Claims" d\textquotesingle identité
lors de la transition entre les comptes, ce qui provoquait parfois des
erreurs de redirection. La résolution de ce problème a nécessité une
configuration minutieuse du middleware
d\textquotesingle authentification et l\textquotesingle application
systématique de l\textquotesingle attribut {[}Authorize(Roles =
"..."){]} sur les action du contrôleur, garantissant ainsi une forme
d'étanchéité entre les différentes interfaces de
l\textquotesingle application.

6 Conclusion

\section{6 Conclusion}\label{conclusion}

Ce travail pratique nous a permis de concevoir une application de
commerce électronique complète, ancrée dans un scénario réel de
restauration moderne. En articulant le projet autour d'ASP.NET Core en
MVC et d'Entity Framework Core, nous avons pu transformer des concepts
théoriques tels que la gestion d\textquotesingle états via les sessions
en une solution fonctionnelle. L\textquotesingle intégration
d\textquotesingle une passerelle de paiement réelle avec Stripe a
constitué une étape déterminante, nous confrontant aux exigences de
sécurité et de synchronisation propres aux transactions financières en
ligne.

Au-delà des simples fonctionnalités de navigation et de commande, le
projet nous a imposé une rigueur de structuration logicielle, notamment
à travers le développement du CartService et la gestion fine des rôles
utilisateurs. Les obstacles rencontrés, qu'il s'agisse de la logique de
conflit entre restaurants ou des défis
d\textquotesingle authentification, ont été des opportunités
d\textquotesingle approfondir notre compréhension des mécanismes du Web.
En conclusion, ce projet représente une immersion concrète dans les
enjeux techniques et commerciaux du commerce électronique, nous offrant
une base solide pour le développement de futures applications web
professionnelles.

\section{}\label{section-10}

\end{document}
